\section{Inline Visual Route: Proof and Formalization}

These figures connect theorem prose, proof sheets, finite witnesses, and formal subsets.


\subsection{Theorem route}

This visual anchor names the objects, review direction, and formula or numeric boundary that keep the figure tied to the argument.


\begin{figure}[htbp]

\centering

\begin{tikzpicture}[>=Latex,node distance=2.4cm]

\node[draw,rounded corners,fill=blue!7,text width=0.28\textwidth,align=center,minimum height=1.0cm] (a) {Claim};

\node[draw,rounded corners,fill=green!7,text width=0.28\textwidth,align=center,minimum height=1.0cm,right=of a] (b) {Proof sheet};

\draw[->,line width=0.8pt] (a) -- node[above,align=center] {assumptions} (b);

\node[draw,rounded corners,fill=orange!8,text width=0.68\textwidth,align=center,below=0.9cm of $(a)!0.5!(b)$] (c) {review question: support class, boundary, falsifier};

\draw[->,dashed] (c.north) -- ($(a)!0.5!(b)$);

\end{tikzpicture}

\caption{Theorem route. Shows the path from public sentence to proof sheet and assumptions. The figure is placed inline in the manuscript; complete machine evidence remains in the evidence package.}

\label{fig:r005-inline-28b_oc133_inline_figures_proof_route-1}

\end{figure}

\subsection{Lean subset route}

This visual anchor names the objects, review direction, and formula or numeric boundary that keep the figure tied to the argument.


\begin{figure}[htbp]

\centering

\begin{tikzpicture}[>=Latex,node distance=2.4cm]

\node[draw,rounded corners,fill=blue!7,text width=0.28\textwidth,align=center,minimum height=1.0cm] (a) {Definition};

\node[draw,rounded corners,fill=green!7,text width=0.28\textwidth,align=center,minimum height=1.0cm,right=of a] (b) {Lean declaration};

\draw[->,line width=0.8pt] (a) -- node[above,align=center] {encoded fragment} (b);

\node[draw,rounded corners,fill=orange!8,text width=0.68\textwidth,align=center,below=0.9cm of $(a)!0.5!(b)$] (c) {review question: support class, boundary, falsifier};

\draw[->,dashed] (c.north) -- ($(a)!0.5!(b)$);

\end{tikzpicture}

\caption{Lean subset route. Shows where mechanization supports a bounded theorem surface. The figure is placed inline in the manuscript; complete machine evidence remains in the evidence package.}

\label{fig:r005-inline-28b_oc133_inline_figures_proof_route-2}

\end{figure}

\subsection{Finite witness route}

This visual anchor names the objects, review direction, and formula or numeric boundary that keep the figure tied to the argument.


\begin{figure}[htbp]

\centering

\begin{tikzpicture}[>=Latex,node distance=2.4cm]

\node[draw,rounded corners,fill=blue!7,text width=0.28\textwidth,align=center,minimum height=1.0cm] (a) {Positive case};

\node[draw,rounded corners,fill=green!7,text width=0.28\textwidth,align=center,minimum height=1.0cm,right=of a] (b) {Negative control};

\draw[->,line width=0.8pt] (a) -- node[above,align=center] {mutation test} (b);

\node[draw,rounded corners,fill=orange!8,text width=0.68\textwidth,align=center,below=0.9cm of $(a)!0.5!(b)$] (c) {review question: support class, boundary, falsifier};

\draw[->,dashed] (c.north) -- ($(a)!0.5!(b)$);

\end{tikzpicture}

\caption{Finite witness route. Shows why negative controls are part of proof readability. The figure is placed inline in the manuscript; complete machine evidence remains in the evidence package.}

\label{fig:r005-inline-28b_oc133_inline_figures_proof_route-3}

\end{figure}

\subsection{Counterexample boundary}

This visual anchor names the objects, review direction, and formula or numeric boundary that keep the figure tied to the argument.


\begin{figure}[htbp]

\centering

\begin{tikzpicture}[>=Latex,node distance=2.4cm]

\node[draw,rounded corners,fill=blue!7,text width=0.28\textwidth,align=center,minimum height=1.0cm] (a) {Assumption set};

\node[draw,rounded corners,fill=green!7,text width=0.28\textwidth,align=center,minimum height=1.0cm,right=of a] (b) {Countermodel};

\draw[->,line width=0.8pt] (a) -- node[above,align=center] {reopening} (b);

\node[draw,rounded corners,fill=orange!8,text width=0.68\textwidth,align=center,below=0.9cm of $(a)!0.5!(b)$] (c) {review question: support class, boundary, falsifier};

\draw[->,dashed] (c.north) -- ($(a)!0.5!(b)$);

\end{tikzpicture}

\caption{Counterexample boundary. Shows how a failed assumption reopens a theorem. The figure is placed inline in the manuscript; complete machine evidence remains in the evidence package.}

\label{fig:r005-inline-28b_oc133_inline_figures_proof_route-4}

\end{figure}

\subsection{Dependency chain}

This visual anchor names the objects, review direction, and formula or numeric boundary that keep the figure tied to the argument.


\begin{figure}[htbp]

\centering

\begin{tikzpicture}[>=Latex,node distance=2.4cm]

\node[draw,rounded corners,fill=blue!7,text width=0.28\textwidth,align=center,minimum height=1.0cm] (a) {Definition};

\node[draw,rounded corners,fill=green!7,text width=0.28\textwidth,align=center,minimum height=1.0cm,right=of a] (b) {Lemma};

\draw[->,line width=0.8pt] (a) -- node[above,align=center] {theorem} (b);

\node[draw,rounded corners,fill=orange!8,text width=0.68\textwidth,align=center,below=0.9cm of $(a)!0.5!(b)$] (c) {review question: support class, boundary, falsifier};

\draw[->,dashed] (c.north) -- ($(a)!0.5!(b)$);

\end{tikzpicture}

\caption{Dependency chain. Shows the proof dependency route as a chain, not a label. The figure is placed inline in the manuscript; complete machine evidence remains in the evidence package.}

\label{fig:r005-inline-28b_oc133_inline_figures_proof_route-5}

\end{figure}

\subsection{Minimality witness}

This visual anchor names the objects, review direction, and formula or numeric boundary that keep the figure tied to the argument.


\begin{figure}[htbp]

\centering

\begin{tikzpicture}[>=Latex,node distance=2.4cm]

\node[draw,rounded corners,fill=blue!7,text width=0.28\textwidth,align=center,minimum height=1.0cm] (a) {Component kept};

\node[draw,rounded corners,fill=green!7,text width=0.28\textwidth,align=center,minimum height=1.0cm,right=of a] (b) {Verdict changes};

\draw[->,line width=0.8pt] (a) -- node[above,align=center] {ablation} (b);

\node[draw,rounded corners,fill=orange!8,text width=0.68\textwidth,align=center,below=0.9cm of $(a)!0.5!(b)$] (c) {review question: support class, boundary, falsifier};

\draw[->,dashed] (c.north) -- ($(a)!0.5!(b)$);

\end{tikzpicture}

\caption{Minimality witness. Shows component independence by keep/drop comparison. The figure is placed inline in the manuscript; complete machine evidence remains in the evidence package.}

\label{fig:r005-inline-28b_oc133_inline_figures_proof_route-6}

\end{figure}
